% !TEX TS-program = pdflatex
% !TEX encoding = UTF-8 Unicode

% This file is a template using the "beamer" package to create slides for a talk or presentation
% - Giving a talk on some subject.
% - The talk is between 15min and 45min long.
% - Style is ornate.

% MODIFIED by Jonathan Kew, 2008-07-06
% The header comments and encoding in this file were modified for inclusion with TeXworks.
% The content is otherwise unchanged from the original distributed with the beamer package.

\documentclass{beamer}
\setbeamercovered{invisible}



% Copyright 2004 by Till Tantau <tantau@users.sourceforge.net>.
%
% In principle, this file can be redistributed and/or modified under
% the terms of the GNU Public License, version 2.
%
% However, this file is supposed to be a template to be modified
% for your own needs. For this reason, if you use this file as a
% template and not specifically distribute it as part of a another
% package/program, I grant the extra permission to freely copy and
% modify this file as you see fit and even to delete this copyright
% notice. 


\mode<presentation>
{
  \usetheme{Warsaw}
  % or ...

  % or whatever (possibly just delete it)
}


\usepackage[english]{babel}
% or whatever

\usepackage[utf8]{inputenc}
% or whatever

\usepackage{times}
\usepackage[T1]{fontenc}
% Or whatever. Note that the encoding and the font should match. If T1
% does not look nice, try deleting the line with the fontenc.

%%% MATH RELATED

\input{macros.tex}

\title{MATH 350-2 Advanced Calculus} 
\subtitle
{} % (optional)

\author[W.R. Casper] % (optional, use only with lots of authors)
{W.R. Casper}
% - Use the \inst{?} command only if the authors have different
%   affiliation.

\institute[California State University Fullerton] % (optional, but mostly needed)
{
  Department of Mathematics\\
  California State University Fullerton}
% - Use the \inst command only if there are several affiliations.
% - Keep it simple, no one is interested in your street address.

\subject{Talks}
% This is only inserted into the PDF information catalog. Can be left
% out. 



% If you have a file called "university-logo-filename.xxx", where xxx
% is a graphic format that can be processed by latex or pdflatex,
% resp., then you can add a logo as follows:

% \pgfdeclareimage[height=0.5cm]{university-logo}{university-logo-filename}
% \logo{\pgfuseimage{university-logo}}



% Delete this, if you do not want the table of contents to pop up at
% the beginning of each subsection:
\AtBeginSubsection[]
{
  \begin{frame}<beamer>{Outline}
    \tableofcontents[currentsection,currentsubsection]
  \end{frame}
}


% If you wish to uncover everything in a step-wise fashion, uncomment
% the following command: 

%\beamerdefaultoverlayspecification{<+->}


\begin{document}

\begin{frame}
  \titlepage
\end{frame}

\begin{frame}{Outline}
  \tableofcontents
  % You might wish to add the option [pausesections]
\end{frame}

% Since this a solution template for a generic talk, very little can
% be said about how it should be structured. However, the talk length
% of between 15min and 45min and the theme suggest that you stick to
% the following rules:  

% - Exactly two or three sections (other than the summary).
% - At *most* three subsections per section.
% - Talk about 30s to 2min per frame. So there should be between about
%   15 and 30 frames, all told.

\section{Real Analysis Lecture 15}


\subsection{Limit of a function}

\begin{frame}{Limit of a function}
Let $(S,d_S)$ and $(T,d_T)$ be metric spaces and $A\subseteq S$.
\pause
Also, let $f: A\rightarrow T$ be a function.
\pause
\begin{defn}
\pause
Then for any accumulation point $p$ of $A$ and element $L\in T$, we say that the \textbf{limit as $x$ approaches $p$ of $f(x)$ is $L$}
\pause
 and write
$$\lim_{x\rightarrow p} f(x) = L$$
\pause
if for all $\epsilon > 0$ there exists $\delta > 0$ such that
$$0 < d_S(x,p) < \delta\ \Rightarrow\ d_T(f(x),L) < \epsilon.$$
\end{defn}
\end{frame}

\begin{frame}{Example}
Prove that
$$\lim_{x\rightarrow 3}\frac{1}{x^2} = \frac{1}{9}$$
\end{frame}

\begin{frame}{Example}
Prove that
$$\lim_{x\rightarrow 4}\frac{\sqrt{x}-2}{x-4} = \frac{1}{4}$$
\end{frame}

\begin{frame}{Limit of a function}
\begin{thm}
\pause
Let $(S,d_S)$ and $(T,d_T)$ be metric spaces, $A\subseteq S$, and $f: A\rightarrow T$ be a function.
\pause
Then
$$\lim_{x\rightarrow a} f(x) = L$$
\pause
if and only if
$$\lim_{n\rightarrow \infty} f(x_n) = L$$
for all sequences $\{x_n\}$ of values in $A\backslash \{a\}$ which converge to $a$.
\end{thm}
\end{frame}

\begin{frame}{Challenge!}
Consider the function
$$f: \mathbb{R}^2\rightarrow \mathbb R$$
between metric spaces with the Euclidean metric defined by
$$f(x,y) = \frac{x^2y}{x^4+y^2}.$$
\pause
\begin{quest}
\pause
What is the limit as $(x,y)\rightarrow (0,0)$?
\end{quest}
\end{frame}

\begin{frame}{Challenge!}
Consider the sequence $\{(x_n,y_n)\}$ defined by
$$(x_n,y_n) = (1/n,0).$$
\pause
Then
$$\lim_{n\rightarrow \infty} (x_n,y_n) = \lim_{n\rightarrow \infty} (1/n,0) = (0,0)$$
\pause
and
$$\lim_{n\rightarrow \infty} f(x_n,y_n) = \lim_{n\rightarrow \infty} 0 = 0.$$
\end{frame}

\begin{frame}{Challenge!}
Consider the sequence $\{(x_n,y_n)\}$ defined by
$$(x_n,y_n) = (0,1/n).$$
\pause
Then
$$\lim_{n\rightarrow \infty} (x_n,y_n) = \lim_{n\rightarrow \infty} (0,1/n) = (0,0)$$
\pause
and
$$\lim_{n\rightarrow \infty} f(x_n,y_n) = \lim_{n\rightarrow \infty} 0 = 0.$$
\end{frame}

\begin{frame}{Challenge!}
Consider the sequence $\{(x_n,y_n)\}$ defined by
$$(x_n,y_n) = (1/n,1/n).$$
\pause
Then
$$\lim_{n\rightarrow \infty} (x_n,y_n) = \lim_{n\rightarrow \infty} (1/n,1/n) = (0,0)$$
\pause
and
$$\lim_{n\rightarrow \infty} f(x_n,y_n) = \lim_{n\rightarrow \infty} \frac{1/n^3}{1/n^4+1/n^2} = \lim_{n\rightarrow \infty} \frac{1/n}{1/n^2+1} = 0.$$
\end{frame}

\begin{frame}{Challenge!}
Consider the sequence $\{(x_n,y_n)\}$ defined by
$$(x_n,y_n) = (1/n,1/n^2).$$
\pause
Then
$$\lim_{n\rightarrow \infty} (x_n,y_n) = \lim_{n\rightarrow \infty} (1/n,1/n^2) = (0,0)$$
\pause
and
$$\lim_{n\rightarrow \infty} f(x_n,y_n) = \lim_{n\rightarrow \infty} \frac{1/n^4}{1/n^4+1/n^4} = \lim_{n\rightarrow \infty} \frac{1}{2} = \frac{1}{2}.$$
\pause
{\color{red}The limit DOES NOT EXIST!}
\end{frame}


\begin{frame}{Limit of a function}
\begin{proof}
$(\Rightarrow)$\quad
Suppose that $\lim_{x\rightarrow a} f(x) =L$ and let $\{x_n\}$ be a sequence of values in $A\backslash\{a\}$ which converges to $a$.\\
\pause
Let $\epsilon > 0$.\\
\pause
Then there exists $\delta > 0$ such that $0 < d_S(x,a) < \delta$ implies $d_T(f(x),L) < \epsilon$.\\
\pause
Also, there exists $N \in\mathbb{Z}_+$ such that $n\geq N$ implies $d_S(x_n,a) < \delta$ for all $n\geq N$.\\
\pause
Therefore since $x_n\neq a$ for all $n$, we have that
$$n\geq N\ \Rightarrow\ 0 < d_S(x_n,a) < \delta\ \Rightarrow\ d_T(f(x_n), L) < \epsilon.$$
\pause
Hence for all $n\geq N$, we know $d_T(f(x_n), L) < \epsilon$.\\
\pause
Since $\epsilon > 0$ was arbitrary, this proves $\lim_{n\rightarrow \infty} f(x_n) = L$.
\end{proof}
\end{frame}

\begin{frame}{Limit of a function}
\begin{proof}
$(\Leftarrow)$\quad
Suppose that for every sequence $\{x_n\}$ of values in $A\backslash\{a\}$ which converges to $a$, we have $\lim_{n\rightarrow \infty} f(x_n) = L$.
\pause
We must show that $\lim_{x\rightarrow a} f(x) = L$.\\
\pause
To see this, we assume otherwise.\\
\pause
Then there exists $\epsilon > 0$ such that for all $\delta > 0$ there exists $x$ with $0 < d_S(x,a) < \delta$ and $d_T(f(x),L) \geq \epsilon$.\\
\pause
For each $n\in \mathbb{Z}_+$, take $x_n\in A$ with $0 < d_S(x_n,a) < 1/n$ and $d_T(f(x_)n,L) \geq \epsilon$.\\
\pause
In this case, $\{x_n\}$ is a sequence of values in $A\backslash\{a\}$ which converges to $a$.\\
\pause
However, $\{f(x_n)\}$ never gets closer than $\epsilon$ to $L$, so it cannot converge to $L$.\\
\pause
This is a contradiction.
\end{proof}
\end{frame}

\subsection{Continuity}

\begin{frame}{Continuity}
Let $(S,d_S)$ and $(T,d_T)$ be metric spaces and $f: S\rightarrow T$ a function.
\pause
\begin{defn}
\pause
Then $f$ is continuous at a point $p\in S$ if 
if for all $\epsilon > 0$ there exists $\delta > 0$ such that
$$d_S(x,p) < \delta\ \Rightarrow\ d_T(f(x),f(p)) < \epsilon.$$
\pause
Equivalently, $f(B_S(p;\delta))\subseteq B_T(f(p);\epsilon)$.
\pause
We say $f$ is continuous on $A\subseteq S$ if $f(x)$ is continuous at every point $p\in A$.
\pause
If $f$ is continuous on $S$, we simply say $f$ is continuous.
\end{defn}
\end{frame}

\begin{frame}{Sequence version}
\pause
\begin{thm}
Let $(S,d_S)$ and $(T,d_T)$ be metric spaces and $f: S\rightarrow T$ a function.
\pause
Then $f$ is continuous at a point $p$ in $S$ if for all sequences $\{x_n\}$ of values in $S$ which converge to $p$, we have
\pause
$$\lim_{n\rightarrow\infty} f(x_n) = f(p).$$
\end{thm}
\pause
Intuitively, continuity means
$$\lim_{n\rightarrow\infty} f(x_n) = f\left(\lim_{n\rightarrow\infty} x_n\right).$$
\end{frame}

\begin{frame}{Compositions}
\begin{thm}
Suppose $(S,d_S)$, $(T,d_T)$, and $(U,d_U)$ are metric spaces and $f: S\rightarrow T$ and $g: T\rightarrow U$ are functions.\\
\pause
If $f$ is continuous at $p$ and $g$ is continuous at $f(p)$, then the composition $h = g\circ f: S\rightarrow U$ defined by $h(x) = g(f(x))$ is continuous at $p$.
\end{thm}
\end{frame}

\begin{frame}{Challenge!}
\begin{prob}
Prove the previous theorem!
\end{prob}
\end{frame}

\begin{frame}{Preimages}
\begin{defn}
Suppose $f: S\rightarrow T$ is a function and that $Y\subseteq T$.\\
\pause
The \textbf{inverse image} or \textbf{preimage} of $Y$ is the set
$$f^{-1}(Y) = \{x: x\in S,\ \text{and}\ f(x)\in Y\}.$$
\end{defn}
\end{frame}

\begin{frame}{Challenge!}
\begin{prob}
Suppose that $S\rightarrow T$ and $Y\subseteq Z\subseteq T$.\\
Prove that
$$f^{-1}(Y)\subseteq f^{-1}(Z).$$
\end{prob}
\end{frame}


\begin{frame}{Open sets}
\begin{thm}
Suppose $(S,d_S)$ and $(T,d_T)$ are metric spaces and $f: S\rightarrow T$ is a function.\\
\pause
Then $f$ is continuous if and only if $f^{-1}(V)$ is open for all open subsets $V\subseteq T$.
\end{thm}
\end{frame}

\begin{frame}{Closed sets}
\begin{thm}
Suppose $(S,d_S)$ and $(T,d_T)$ are metric spaces and $f: S\rightarrow T$ is a function.\\
\pause
Then $f$ is continuous if and only if $f^{-1}(A)$ is closed for all closed subsets $A\subseteq T$.
\end{thm}
\end{frame}



\end{document}


